% ------------------------------------------------------------------------
% file `v1-exercise-6-solution-body.tex'
%   in folder `xsim/'
%
%     solution of type `exercise' with id `6'
%
% generated by the `solution' environment of the
%   `xsim' package v0.21 (2022/02/12)
% from source `v1' on 2022/09/08 on line 281
% ------------------------------------------------------------------------
    \begin{enumerate}
        \item 0,5 point pour $m$.\\
              0,5 point pour $p$.\\
              0,5 point si l'élève cherche la racine.\\
              0,5 point pour la racine correcte.\\
              -0,5 point si la réponse n'est pas formulée sous forme d'une phrase.

              $f_A(x) = -\dfrac{1}{2}x + 6$
              \begin{multicols}{2}
                  \textbf{$\bm{m}$ ?}
                  \medbreak
                  $m=\dfrac{-1}{2}=-\dfrac{1}{2}$
                  \columnbreak

                  \textbf{$\bm{p}$ ?}
                  \medbreak
                  $p=6$
              \end{multicols}
              Après combien de jours la hauteur sera-t-elle nulle ?
              \begin{align*}
                  0  & = -\dfrac{1}{2}x + 6 \\
                  -6 & = -\dfrac{1}{2}x     \\
                  12 & = x
              \end{align*}
              L'évaporation de l'huile contenue dans l'éprouvette A sera complète après 12 jours.
        \item 0,5 point pour $m$.\\
              0,5 point pour $p$.\\
              0,5 point si l'élève cherche la racine.\\
              0,5 point pour la racine correcte.\\
              -0,5 point si la réponse n'est pas formulée sous forme d'une phrase.

              $f_B(x)=-\dfrac{1}{4}x + 4$
              \begin{multicols}{2}
                  \textbf{$\bm{m}$ ?}
                  \medbreak
                  $m=\dfrac{-1}{4}=-\dfrac{1}{4}$
                  \columnbreak

                  \textbf{$\bm{p}$ ?}
                  \medbreak
                  $p=4$
              \end{multicols}
              Que est l'abscisse du point d'intersection de $f_A$ et $f_B$ ?
              \begin{align*}
                  -\dfrac{1}{2}x + 6             & = -\dfrac{1}{4}x + 4 \\
                  -\dfrac{1}{2}x + \dfrac{1}{4}x & = 4 - 6              \\
                  -\dfrac{2x}{4} + \dfrac{x}{4}  & = -2                 \\
                  \dfrac{-x}{4}                  & = -2                 \\
                  -x                             & = -8                 \\
                  x                              & = 8
              \end{align*}
              Les quantités d'huile contenues dans les éprouvettes seront égales après 8 jours.
    \end{enumerate}
